\documentclass[a4paper,onecolumn,11pt]{article}
\usepackage{amsfonts}
\usepackage{amssymb}
\usepackage{amsmath}
\usepackage{latexsym}
\usepackage{cuted}
\usepackage{graphicx}
\usepackage{biblatex}
\usepackage[skip=2pt, font=scriptsize]{caption}
\usepackage{indentfirst}
\usepackage{enumitem}
\hyphenpenalty1000
\addbibresource{main.bib}

%-----------------------------------------------------------------------------
\usepackage{xcolor}
\usepackage{listings}
\definecolor{codegreen}{rgb}{0,0.6,0}
\definecolor{codegray}{rgb}{0.5,0.5,0.5}
\definecolor{codepurple}{rgb}{0.58,0,0.82}
\definecolor{backcolour}{rgb}{0.95,0.95,0.92}
\lstdefinestyle{mystyle}{
    backgroundcolor=\color{backcolour},   
    commentstyle=\color{codegreen},
    keywordstyle=\color{magenta},
    numberstyle=\tiny\color{codegray},
    stringstyle=\color{codepurple},
    basicstyle=\ttfamily\scriptsize,
    breakatwhitespace=false,         
    breaklines=true,                 
    captionpos=b,                    
    keepspaces=true,                 
    numbers=left,                    
    numbersep=5pt,                  
    showspaces=false,                
    showstringspaces=false,
    showtabs=false,                  
    tabsize=2
}
\lstset{style=mystyle}
%------------------------------------------------------------------------------

\newcounter{vatn} % Assumption counter
\setcounter{vatn}{1} % Set to 1

%------------------------------------------------------------------------------

\begin{document}

\title{Inverse Kinematics and Workspace Analysis of X-Y Pedestal
with Joint and Pointing Misalignment } 
\author{Buğra Coşkun\thanks{E-mail:
\texttt{bugra.coskun@profen.com}}}

\date{\vspace{-2ex}\scriptsize05.09.2023}
\maketitle
\vspace{-5ex}
\normalsize

\begin{abstract}
    We present an inverse kinematics solution for X-Y pedestals with joints and
    pointing misalignment. We use vector analysis for optimal computational
    complexity of the inverse kinematics solution. We show how misalignment in
    the joints affect the workspace of the pedestal and present examples in
    Octave and C for the application of the solution. 
\end{abstract}

\normalsize
\section*{Introduction}
The 2R serial manipulator configurations are popular solutions for satellite
ground stations with mechanical tracking capabilities.
Altitude-azimuth pedestals are widely used for their relative simplicity and
cost-effectiveness but they proove less than optimal for high-elevation 
applications because
of the zenith problem creating an increase in needed azimuth
velocities\cite{ji2003} as the tracked target gets closer to the azimuth joint
axis. To combat the problem with tracking objects in high
elavations the X-Y configuration is another popular option for ground
terminals. In this configuration the joints are moving slower compared to
altitude-azimuth configuration at higher elevations but a similar mechanical
limitation arises that decreases the manipulator's effectiveness at relatively
low elevations. As the need for low elevation tracking is quite rare the X-Y
configuration usually happens to be an optimal solution for the pedestal
choice\cite{willey2000}. Research into the kinematics and dynamics for the X-Y
manipulator has been previously explored using D-H parameters for the inverse
kinematics problem \cite{tehrani2014}. The synthesized solutions use
trigonometric functions but the added misalignment properties increases the
complexity of this method. We shall instead employ vector analysis approach to
keep the complexity of the equations relatively tame.

\normalsize
\section{Forward Kinematics}
For the forward kinematics we first consider the rotation of the Y-joint,
because the Y-joint axis is transformed by the X-joint rotation, considering
this rotation first will preserve us from calculating the new axis of the
Y-joint. This is an example of extrinsic rotation. The pedestal orientation
vector, Y-joint axis and X-joint axis are denoted as $\hat p_0, \hat w_0, \hat
x_0$ respectively.
\begin{equation}
    \hat w_0 = \begin{bmatrix}
            \cos(\lambda) \\
            \sin(\lambda) \\
            0
        \end{bmatrix} , \;
    W_0 = \begin{bmatrix}
        0 & -k_3 & k_2 \\
        k_3 & 0 & -k_1 \\
        -k_2 & k_1 & 0
    \end{bmatrix}
    \; \therefore \; W_0 \vec v = \hat w_0 \times \vec v.
\end{equation}
The rotation matrix in $SO(3)$ group around an arbitrary axis is given by the 
Rodrigues' rotation formula:
\begin{equation}
    R_{\hat w_0}(\phi) = I + (\sin \phi) W_0 + (1 - \cos \phi) W_0^2,
\end{equation}
The standard rotation matrix around the $\hat x$ axis is in the form:
\begin{equation}
    R_x ( \psi ) = \begin{bmatrix}
        1 & 0 & 0 \\
        0 & \cos \psi & -\sin \psi \\
        0 & \sin \psi & \cos \psi
    \end{bmatrix},
\end{equation}
Now that we have the rotation matrices that we require, the forward kinematic
equation becomes:
\begin{equation}
    \hat l = R_x ( \psi ) R_{w_0} ( \phi ) \hat p_0,
\end{equation}
The resulting line-of-sight vector is in the pedestal frame, to get the
orientation in the NWU(north-west-up) frame we matrix multiply the
LOS vector with the the inverse of the NWU rotation matrix defined in the 
pedestal basis. For further details on the NWU matrix and frame see Appendix B.
\begin{equation}
    \hat l_N = R_N^{-1} \hat l.
\end{equation}
The resulting azimuth and elevation degrees is given by the cartesian to
spherical conversion. For our applications we do not need the radial distance
and it is more convenient to use the elevational angle instead of the polar
angle. We define the function and it's inverse for these conversions as such:
\begin{equation}
    \Lambda : \mathbb{R}^2 \mapsto \mathbb{R}^3, \; 
    \phi \in (-\pi, \pi],\; 
    \theta \in \left[-\frac{\pi}{2},\frac{\pi}{2}\right],
\end{equation}
\begin{equation}
    \Lambda
    \left( \begin{bmatrix} \phi \\ \theta \end{bmatrix} \right)
    =
    \begin{bmatrix}
        \cos( \theta ) \cos(\phi) \\
        \cos( \theta ) \sin(\phi) \\
        \sin ( \theta )
    \end{bmatrix}
    \text{ and }
    \Lambda^{-1}( \hat v ) = \begin{bmatrix}
        \text{atan2}(v_2, v_1)
        \\
        \sin ^{-1} ( v_3 )
        \end{bmatrix},
\end{equation}
\begin{equation}
    \therefore \begin{bmatrix} - \alpha \\ \beta \end{bmatrix} =
        \Lambda^{-1}(\vec l_N).
\end{equation}
Where $\alpha$ and $\beta$ are the azimuth and elevation degrees of the
resulting pedestal orientation. The azimuth angle is taken the negative of
because azimuth degree is CW positive contrary to the conventional spherical
coordinates. 

\section{Inverse Kinematics}
For a pedestal orientation we are most often going to be supplied with an
azimuth and elevation angles, thus we first calculate the orientation vector
using the function from Equation (7).
\begin{equation}
    \hat l = \Lambda\left(\begin{bmatrix} - \alpha \\ \beta \end{bmatrix}\right) =
    \begin{bmatrix}
        l_{1} \\
        l_{2} \\
        l_{3}
    \end{bmatrix}
\end{equation}
For the inverse kinematics we need to know what Y-joint angle brings the
target vector into a position that will reach the target orientation when
subjected to a rotation about the X-joint.

\begin{equation}
    \mathbb{V} = \{ \vec v \in \mathbb{R}^3 : || \vec v || = 1 \}
\end{equation}
\begin{equation}
    c \in [-1,1] \rightarrow \mathbb{S}_{\hat w}(c) = \{ \hat v :
        \hat v \cdot \hat w = c \}.
\end{equation}
\begin{equation}
    \forall \hat v \in \mathbb{S}_{\hat w}(c), \forall \theta \in (-\pi,\pi]
        \rightarrow R_{\hat w}(\theta)\hat v \in \mathbb{S}_{\hat w}(c)
        \because R_{\hat w}(\theta) \in SO(3).
\end{equation}
\begin{equation}
    \hat l \in \mathbb{S}_{\hat x} (\hat l \cdot \hat x) \therefore 
    R_{\hat x} (\psi) R_{\hat w_0}(\phi) \hat p_0 = \hat l
    \Leftrightarrow R_{\hat w_0}(\phi) \hat p_0 
    \in \mathbb{S}_{\hat x}(\hat l \cdot \hat x)
\end{equation}


\section*{Acknowledgements}
The author is thankful to Mustafa Çelik, Cem Doğan and Kadir Koca for their 
guidance on X-Y pedestals as well as to Merve Güvenç for the inspiring 
conversations and tips for structuring this paper.


\printbibliography

\section*{Appendix A: North-West-Up Frame}
\section*{Appendix B: Octave Code - Inverse Kinematics}
\lstinputlisting[language=Matlab]{xyink.m}
\end{document}
